\subsection{滤子的定义}

定义 1. 集合 X 上的滤子是 X 的子集所成的一个集合 F，它具有下列性质：

（F \(_{I}\)）X 的每个包含 F 中某个集合的子集都属于 F。

\((\mathbf{F}_{\mathrm{II}})\) \(\mathfrak{F}\) 的集合的每个有限交都属于 \(\mathfrak{F}\)。

\((\mathbf{F}_{\mathrm{III}})\) 空集不属于 \(\mathfrak{F}\)。

由 \((\mathbf{F}_{\mathbf{II}})\) 与 \((\mathbf{F}_{\mathbf{III}})\) 可知，F 的集合的每个有限交都非空。

X 上的滤子 \(\mathfrak{F}\) 在 X 上定义了一个结构，其公理为 \((\mathrm{F_I})\)、\((\mathrm{F_{II}})\) 与 \((\mathrm{F_{III}})\)；这个结构称为滤集结构，而赋以该结构的集合 X 称为被 \(\mathfrak{F}\) 滤化的集合。

公理 \((\mathbf{F}_{\mathrm{II}})\) 等价于下列两条公理的合取：\((\mathbf{F}_{\mathrm{II}a})\) F 中两个集合的交属于 F。

\((\mathbf{F}_{\mathrm{II}b})\) X 属于 \(\mathfrak{F}\)。

公理 \((\mathbf{F}_{\mathbf{II}\ \mathbf{b}})\) 与 \((\mathbf{F}_{\mathbf{III}})\) 表明空集上不存在滤子。

为使满足 \((\mathbf{F}_{\mathbf{I}})\) 的子集族也满足 \((\mathbf{F}_{\mathbf{II}\mathbf{b}})\)，必须且只需它非空。满足 \((\mathbf{F}_{\mathbf{I}})\) 的子集族也满足 \((\mathbf{F}_{\mathbf{III}})\)，当且仅当它异于 \(\mathfrak{P}(\mathbf{X})\)。

滤子的例子. 1) 若 \(\mathbf{X} \neq \emptyset\)，则仅由 \(\mathbf{X}\) 组成的子集族是 \(\mathbf{X}\) 上的一个滤子。更一般地，\(\mathbf{X}\) 中包含给定非空子集 \(\mathbf{A}\)（\(\mathbf{X}\) 的）的所有子集所成之集是 \(\mathbf{X}\) 上的一个滤子。

\begin{enumerate}
\def\labelenumi{\arabic{enumi})}
\setcounter{enumi}{1}
\item
  在拓扑空间 X 中，X 的任一非空子集 A 的所有邻域所成之集（特别地，X 中一点的所有邻域所成之集）是一个滤子，称为 A 的邻域滤子。
\item
  若 X 是无限集，则 X 的有限子集的余集组成一个滤子的元素。由 \(\geqslant0\) 的整数组成的集 N 的有限子集的余集所成的滤子称为 Fréchet 滤子。
\end{enumerate}

